\documentclass[11pt]{article}
\usepackage[margin=1in]{geometry}
\usepackage{amsmath,amssymb,amsfonts,mathtools}
\usepackage{array,booktabs,enumitem,graphicx,xcolor,hyperref}
\usepackage[T1]{fontenc}
\usepackage[utf8]{inputenc}
\title{Hamiltonian construction of W-gravity actions}
\author{A. Mikovic}
\date{}
\begin{document}
\maketitle
\begin{abstract}

We show that all W-gravity actions can be easilly constructed and understood from the point of view of the Hamiltonian formalism for the constrained systems. This formalism also gives a method of constructing gauge invariant actions for arbitrary conformally extended algebras.

\end{abstract}

\section{Introduction}

A B C D E F G H I J K L M N O P Q R S T U V W X Y Z `@=11 \#1 \textbackslash{}@@underline\#1 math expression `\textbackslash{}'= ' \textasciicircum{} @s \#1 \#1 \textbackslash{}@m 1200 1440 1728 2074 2488 2986 =cmmi10\textbackslash{}@magscale5 ='177 =cmsy10\textbackslash{}@magscale5 ='60 =lasy10\textbackslash{}@magscale5 =cmr10\textbackslash{}@magscale6 @ @ @ \textbackslash{}@ne \textbackslash{}@ne \textbackslash{}@ne @ @ @ @@ @@ @@

\section{Method}

L L \#1 \textasciicircum{} \#1 \#1 \_ \#1 \#1 /\#1 \#1 \#1 \#1 \#1 \#1 \#1 \#1 \#1 \#1 \#1 \#1 | \#1 | \#1 \#1 \#1 \#1\#2 \#1 |\#2 \#1 | \#1 | \#1 \#1 | \#1 | \#1 \#1 \#1 \#1 | \#1 |

\begin{equation}
\label{eq:compact}
L(\theta) = \sum_{i=1}^{n} \ell(x_i, \theta)
\end{equation}
Equation~\ref{eq:compact} is included to keep the sample paper-like while remaining small.

\section{Conclusion}

\#1\#2 1 .5ex 1 .4ex \#1\#2 1 .5ex 1 .3ex \#1\#2 1 .5ex \#2 \#1\#2 d \#1 d \#2 \#1\#2 \#1 \#2 \#1\#2 \#1 | \#2 \#1\#2\#3 \textasciicircum{}2 \#1 \#2 \#3 \#1\#2 \textasciicircum{} \#1 \#1\#2 \#1\#2 \_ \#1 \#1 \#1 \#1
\end{document}
